\documentclass[a4paper,12pt]{article}
\setlength{\headheight}{68.803pt}
\addtolength{\topmargin}{-56.803pt}

\usepackage[utf8]{inputenc}
\usepackage[T1]{fontenc}
\usepackage{lmodern}
\usepackage[nopatch=footnote]{microtype}
\usepackage{fvextra}
\DefineVerbatimEnvironment{Verbatim}{Verbatim}{
  breaklines=true,
  breakanywhere=true,
  frame=single,
  fontsize=\small
}

\usepackage[
  left=1in,
  right=1in,
  top=3cm,
  bottom=1in
]{geometry}

\usepackage[
    colorlinks=true,
    linkcolor=black,
    urlcolor=blue,
    citecolor=blue,
    pdfborder={0 0 0},
    pdfauthor={Massimo Mantineo},
    pdftitle={Condition Variables e Sincronizzazione Avanzata},
    pdfsubject={Laboratorio di Reti e Sistemi Distribuiti: HANDS-ON 13},
]{hyperref}

\usepackage{graphicx}
\graphicspath{{../../assets/}}
\usepackage{tikz}
\usetikzlibrary{positioning, arrows.meta, shapes.symbols}
\usepackage{listings}
\usepackage{xcolor}
\usepackage{amsmath, amssymb}
\usepackage{fancyhdr}
\usepackage{setspace}
\usepackage{pgfplots}
\pgfplotsset{compat=1.17}
\usepackage{booktabs}
\usepackage[skins, listings]{tcolorbox}
\usepackage{float}

\newtcblisting{terminal}[1][]{
    listing only,
    colback=black!85!white,
    colframe=black!70!white,
    colupper=green!80!white,
    coltitle=white,
    title=#1,
    fonttitle=\bfseries\sffamily,
    listing options={
        language={},
        basicstyle=\ttfamily\small\color{green!80!white},
        backgroundcolor={},
        frame=none,
        breaklines=true,
        showstringspaces=false,
        numbers=none
    },
    enhanced,
    drop shadow,
    arc=3mm
}

\newtcblisting{ccode}[1][]{
    listing only,
    colback=blue!3!white,
    colframe=blue!70!black,
    title=#1,
    fonttitle=\bfseries\sffamily,
    listing options={
        style=cstyle,
        backgroundcolor={},
        frame=none
    },
    enhanced,
    drop shadow,
    arc=3mm
}

\setlength{\footskip}{50pt}

\lstdefinestyle{cstyle}{
    language=C,
    basicstyle=\ttfamily\small,
    numbers=left,
    numberstyle=\tiny,
    stepnumber=1,
    numbersep=5pt,
    backgroundcolor=\color{white},
    showspaces=false,
    showstringspaces=false,
    showtabs=false,
    frame=single,
    tabsize=4,
    captionpos=b,
    breaklines=true,
    breakatwhitespace=true,
    keywordstyle=\color{blue}\bfseries,
    commentstyle=\color{green!50!black},
    stringstyle=\color{red},
}
\lstset{style=cstyle}

\pagestyle{fancy}
\fancyhf{}

\fancyhead[L]{%
  \parbox[b]{0.45\textwidth}{%
    \small
    Student Name: Massimo Mantineo\\
    Student ID: 541924
    \vspace{20pt}
  }%
}
\fancyhead[C]{%
  \parbox[b]{0.2\textwidth}{%
    \centering
    \normalsize \bfseries
    LabReSiD25\\
    Hands-On 13
    \vspace{20pt}
  }%
}
\fancyhead[R]{%
  \parbox[b]{0.35\textwidth}{%
    \flushright
    \includegraphics[width=3.5cm]{\detokenize{logo_unime_orizontale.png}}
    \vspace{15pt}
  }%
}

\renewcommand{\contentsname}{Indice}

\title{\vspace{-2cm}Relazione: Condition Variables e Sincronizzazione Avanzata\\
\large (Hands-On 13)}
\author{Massimo Mantineo \\ Matricola: 541924}
\date{MESSINA, A.A. 2024/2025}

\begin{document}

\begin{titlepage}
    \thispagestyle{empty}
    \centering
    \vspace*{1cm}
    \includegraphics[width=6cm]{\detokenize{logo_unime_orizontale.png}}\par\vspace{1cm}
    \vspace{2cm}
    {\Large \textbf{Università degli Studi di Messina}\par}
    {\large Dipartimento MIFT - Corso di Laurea Triennale in Informatica (L-31)\par}
    \vspace*{1cm}
    {\large Laboratorio di Reti e Sistemi Distribuiti\par}
    \vspace{1cm}
    {\Large \textbf{HANDS-ON 13:}\\
    Condition Variables, Wait/Signal e Barriere\par}
    \vspace{2cm}
    {\large Massimo Mantineo \par}
    {\large Matricola: 541924 \par}
    \vfill
    {\large Messina, A.A. 2024/2025 \par}
\end{titlepage}

\clearpage
\pagestyle{fancy}
\tableofcontents
\newpage

\section{Problem Statement}
L'Hands-On 13 chiude il cerchio e il quadro dei costrutti di sincronizzazione introducendo le **Condition Variables** e i meccanismi signal-driven. 
Lo scopo dell'esercitazione richiede la loro applicazione pratica lungo due scenari:
1. \textbf{Simulazione di una Thread Join}: Costruire una logica padre-figlio impiegando segnali di terminazione espliciti, interdicendo esplicitamente l'operazione \texttt{pthread\_join()} e ricorrendo alle chiamate \texttt{wait()} e \texttt{signal()} per mettere in attesa il server master al compiersi dell'esatta direttiva asincrona dello slave.
2. \textbf{Barriera di Sincronizzazione con Broadcast}: Ricreare l'impianto della barriera (struttura vista all'Hands-on 12) rimuovendo il ricorso a cascatismi di \textit{sem\_post} su semafori ed integrando al suo posto la logica scalabile delle funzioni \texttt{broadcast()} delle variabili condizionali.

\section{Stato dell'Arte}
Le Variable Condition segnano un passaggio dai lock puramente "esclusivi" (Mutex) e limitativi (Semafori) ad una forma di cooperazione logica di alto livello.
\begin{itemize}
    \item \textbf{Condition Variable (\texttt{pthread\_cond\_t})}: Astrazione asincrona in C/POSIX che gestisce implicitamente una "Waiting room" (coda dei sospesi). A differenza dei semafori, queste variabili non possiedono una logica interiore (come un count numerico), motivo per il quale non possono **mai** esser utilizzate senza l'accoppiamento stretto, e rigoroso, con un \textit{Mutex} di riferimento per proteggere la booleana/stato del costrutto in memoria che decreta l'attesa.
    \item \textbf{Wait e Signal}: \texttt{pthread\_cond\_wait} è in grado atomicamente di sbloccare il mutex ed entrare a basso pre-emption a fare nulla fin quando non viene scossa da remoto. Il "remote shacker" da un altro thread chiamerà la corrispondente \texttt{pthread\_cond\_signal}.
\end{itemize}

\section{Metodologia ed Implementazione}

\subsection{Simulazione di un Join (\texttt{ex1\_join.c})}
Questa tipologia di codice mima al 100\% la gestione in background dello scheduler per attendere flussi operativi prima di uscire in un \texttt{return 0;}.
\begin{ccode}[Snippet Esercizio 1: Padre vs Figlio]
// T-PADRE: Si blocca se non pronto
pthread_mutex_lock(&lock);
while (child_done == 0) {  // Uso un while ciclico di protezione base
    pthread_cond_wait(&cond, &lock);
}
// Se e' qui in basso, cond\_signal e' arrivato e lock e' stato ripreso in mano!
pthread_mutex_unlock(&lock);

// T-FIGLIO: Avvisa la terminazione
pthread_mutex_lock(&lock);
child_done = 1;              
pthread_cond_signal(&cond); // Rende udibile ai pendenti il cambiamento
pthread_mutex_unlock(&lock);
\end{ccode}

\subsection{Barriera a Broadcast (\texttt{ex2\_barrier.c})}
Prevede la sostituzione radicale nel \textit{modus operandi} di blocco della barriera implementando la \textit{funzione massiva POSIX d'ambiente}: cond\_broadcast.
\begin{ccode}[Snippet Esercizio 2: Arrivo globale al punto zero]
if (count == N) {
    // Ultimo thread arrivato al checkpoint
    pthread_cond_broadcast(&cond); // -> Sveglia massivamente da 0 N thread in coda
} else {
    // Si accodano attendendo
    while (count < N) {
        pthread_cond_wait(&cond, &lock);
    }
}
\end{ccode}


\subsection{Istruzioni di Esecuzione}
L'applicativo viene gestito interamente tramite Makefile:
\begin{terminal}[Build & Run]
$ make
$ ./programma
\end{terminal}

\section{Risultati}
Le compilazioni e le catture da log sul terminale attestano il soddisfacimento dei vincoli imposti al design.

\subsection{Output Simulazione Join}
\begin{terminal}[System V Message Queues - Risultati]
massimo@PORTATILE-Linux:~/Scrivania/LabReSid24-25/hands-on/ho13$ ./ex1_join
Padre: Creo il thread figlio.
Padre: In attesa della terminazione del figlio...
Thread Figlio: Iniziato. Faccio un po' di lavoro...
Thread Figlio: Lavoro terminato. Avviso il padre...
Padre: Il figlio ha confermato la terminazione (Join simulato con successo).
\end{terminal}
La dimostrazione concreta che il Thread chiamante non incorre ad uscita forzata nonostante la mancanza di una vera \texttt{pthread\_join}. Questo garantisce lo shutdown dolce.

\subsection{Output Barriera a Variabile Condizionale}
\begin{terminal}[POSIX Message Queues - Risultati]
Thread 0: Arrivato alla barriera. (Presenti: 1/4)
Thread 1: Arrivato alla barriera. (Presenti: 2/4)
Thread 2: Arrivato alla barriera. (Presenti: 3/4)
Thread 3: Arrivato alla barriera. (Presenti: 4/4)
Thread 3: Sono l'ultimo arrivato! Sblocco tutti quanti!
Thread 3: Barriera superata! Avvio simultaneo.
Thread 1: Barriera superata! Avvio simultaneo.
Thread 0: Barriera superata! Avvio simultaneo.
Thread 2: Barriera superata! Avvio simultaneo.
\end{terminal}
Il trace del Thread 3 fa scattare atomicamente in concorrenza l'operazione successiva al blocco condizionale rendendo perfetta nel \textit{timestamp} la coordinazione d'avvio.

\section{Conclusioni e Sviluppi Futuri}
Le Condition Variables sono apparse e si sono dimostrate infinitamente più pulite per le gestioni a coda massiva. Il Broadcast ne esalta a pieno la supremazia su semafori quando quest'ultimi sarebbero costretti a gestire la fine scalata (loop for di $N$ \texttt{sem\_post}) o ad ingolfare l'handler chiamante limitando fortemente l'allargabilità orizzontale. 

\noindent \textbf{Sviluppi futuri:}
\begin{itemize}
    \item \textbf{Spurious Wakeups}: Aggiunta di mitigazioni estensive per proteggere i contatori dall'insorgere occasionale di finti risvegli della CPU che attivano chiamate signal del tutto impreviste ma documentate sui layer UNIX POSIX.
    \item \textbf{High-Level Monitor}: Sviluppo di una architettura a "Monitor", incapsulando tutte le condition variables e mutex per il design pattern Prod-Cons mediante l'utilizzo dell'Incapsulamento Object-Oriented assemblandosi sui framework \texttt{C++11/C++20 std::condition\_variable}.
\end{itemize}

\end{document}
